\silentchapter{Proposisjonar}{proposisjonar}

\prop{1}{}{Formverda er alt som er tilfelle.}
\prop{1.1}{o}{Formverda er totaliteten av realiserte konfigurasjonar, ikkje av ting\footnote{Den underliggande ontologien skil ikkje mellom ei fysisk samanføyning, eit distribuert nettverk av mikroteenester eller eit sekvensielt endestadium i eit systemkall.}.}
\prop{1.11}{t}{Alt som har form, har nettopp denne forma og inga anna\footnote{Ein maskintilstand i eit lukka system fastslår identiteten sin absolutt i augneblikket av eksekvering.}.}
\prop{1.12}{t}{At ein konfigurasjon tek éi eksakt form, der uendeleg mange andre var logisk moglege, krev ei årsaksforklaring utover konfigurasjonen sjølv\footnote{Ei spesifikk minneallokering refererer naudsynleg utover sine eigne byte til prosessen som gav kommandoen om ho.}.}
\prop{1.121}{t}{Forklaringa ligg i relasjonen mellom den manifesterte konfigurasjonen og totaliteten av formene som var moglege, men urealiserte\footnote{Ei grensesnittutforming får tyngda si ikkje frå piksla åleine, men frå frastøytinga av alle dei latente interaksjonsstiane brukaren vart forhindra frå å ta.}.}
\prop{1.1211}{t}{Dette er ein relasjonell, ikkje absolutt, eigenskap\footnote{Informasjonsverdien til eit signal er eksklusivt gitt ved rommet av signal som kunne ha vore sende, men som vart utelatne.}.}
\prop{1.2}{d}{Eit \textbf{objekt}\footnote{I ein databasemodell er objektet den atomære oppføringa; i ein asynkron arkitektur er det den udelelege pakka av tilstandsendring.} er den enklaste bestanddelen i ei form.}
\prop{1.21}{t}{Objektet er udeleleg og uforanderleg\footnote{Dekomponering nedanfor dette abstraksjonsnivået oppløyser funksjonen og sprengjer den ontologiske ramma for handlingsrommet.}.}
\prop{1.211}{t}{Objekta utgjer substansen i formverda\footnote{Repertoaret av primitivar definerer det ytre sjiktet for kva som i det heile let seg kompilere.}.}
\prop{1.2111}{t}{Sidan objektnaturen ber i seg moglegheita til å inngå i konfigurasjonar, er alle moglege former alt gjevne med objekta\footnote{Ein kryptografisk protokoll har allereie skissert topologien for alle framtidige transaksjonar lenge før koden er køyrd for første gong.}.}
\prop{1.21111}{o}{Moglegheita er logisk til stades fordi objekta er det, ikkje fordi nokon tenkjer ho\footnote{Ein parser evaluerer syntakstreet fullstendig blindt for intensjonen til programmeraren.}.}
\prop{1.22}{d}{Ein \textbf{konfigurasjon}\footnote{Grafen av nodar og kantar bestemmer eigenskapane til systemet strengt utan omsyn til kva domene nodane representerer.} er ei bestemd samansetjing av objekt.}
\prop{1.221}{t}{Forma er strukturen til konfigurasjonen\footnote{Det materielle substratet kan skiftast frå silisium til karbon, men overføringsfunksjonen står urørd.}.}
\prop{1.3}{d}{\textbf{Formrommet}\footnote{Søkerommet for ein genetisk algoritme er isomorft med alle morfologiske tilstandar klassen tillét.} (\textit{morphospace}) til ein klasse er mengda av alle moglege konfigurasjonar av objekt i denne klassen.\footnote{Raup (1966); Mitteroecker \& Huttegger (2009)}}
\prop{1.31}{d}{Kvart realisert objekt er eitt eksakt punkt i dette n-dimensjonale rommet\footnote{Tilstandsvektoren i ein kognitiv arkitektur tilordnar konfigurasjonen ein irreversibel og eintydig indeks i rommet.}.}
\prop{1.32}{d}{Det empiriske formrommet er alltid ein projeksjon\footnote{Kompresjon av høgdimensjonal brukaråtferd ned til skalerbare metrikkar tvingar fram eit tap av oppløysing i faseovergangen.}.}
\prop{1.321}{t}{Inga endeleg mengd parametrar fangar den latente kompleksiteten heilt og fullt.\footnote{Thompson (1917)}}
\prop{1.3211}{o}{Kvar projeksjon er eit val: ho avdekkjer somme relasjonar og døymer andre\footnote{Kvar funksjonell optimalisering av systemarkitekturen inneber ein samstundes degenerering av andre, unytta kapabilitetar.}.}
\prop{1.32111}{o}{Det finst ingen nøytral projeksjon\footnote{Kvar aggregering av loggdata brenn bort informasjon; det som gjenstår er forma av eit vilkårleg føremål.}.}
\prop{1.32112}{o}{Projeksjonen er naudsynleg for empirisk handtering, men ho konstituerer ikkje forma\footnote{Strukturen til ei nettside som teiknast på skjermen er projeksjonen; trestrukturen i minnet ber den logiske forma.}.}
\prop{1.321121}{o}{Forma eksisterer uavhengig av målinga\footnote{Tilstanden kviler uavhengig av observasjonsmekanismen.}.}
\prop{1.33}{d}{Klassen er den funksjonelle kategorien som avgrensar formrommet\footnote{Typing set ramma for operasjonane; det som ikkje kan validerast innanfor typesystemet, fell utanfor det moglege rommet.}.}
\prop{1.331}{t}{Ho etablerer grensene; seleksjonstrykka avgjer fordelinga innanfor grensene\footnote{Nettverkets bandbreidde og den menneskelege merksemda agerer som filtermekanismar for kva tilstandar som akkumulerast og overlever i systemet.}.}
\prop{1.4}{d}{Formrommet deler seg topologisk i tre regionar: dei busette, dei opne og dei forbodne\footnote{Designprosessen er redusert til styring av grensesnitta mellom desse tre ontologiske fasane.}.}
\prop{1.41}{o}{Dei busette regionane utgjer det historiske arkivet\footnote{Ein versjonskontrollert logg inneheld heile slektstreet av levedyktige mutasjonar over tid.}.}
\prop{1.411}{o}{Dei fungerer som ankerpunkt for all framtidig navigasjon\footnote{Grensesnittkonvensjonar som dra-og-slepp herdar til faste stasjonar systemet lyt passere for å sikre framkommelegheit.}.}
\prop{1.4111}{t}{Tyngda av eit ankerpunkt er proporsjonal med okkupasjonstida og talet på uavhengige tradisjonar som har konvergert dit\footnote{Adopsjonen av eit standardisert protokollformat i distribuerte nettverk genererer eit gravitasjonsfelt som tvingar nye nodar inn i same bane.}.}
\prop{1.41111}{t}{Konvergens frå uavhengige kjelder vitnar sterkare om eit reelt optimum i landskapet enn konvergens frå éin einskild tradisjon\footnote{Når urelaterte maskinlæringsmodellar oppdagar same heuristikk, speglar dette ein djupare lovmessigheit i datasettet, ikkje ein artefakt ved treninga.}.}
\prop{1.42}{o}{Dei opne regionane utgjer det \textit{tilstøytande moglege}\footnote{Det tilstøytande moglege utgjer den presise radiusen av innovasjon som kodesystemet toler ved neste iterasjon.}.}
\prop{1.421}{o}{Dei er tilgjengelege, men uaktualiserte\footnote{Avkutta kodespor kviler i basen, kompilert, men uoppkalla fram til eit eksternt handlingsmønster bryt straumen.}.}
\prop{1.4211}{t}{Det tilstøytande moglege er ikkje uendeleg\footnote{Det kombinatoriske rommet av lovlege operasjonar krympar asymmetrisk jo tettare integrert systemet er.}.}
\prop{1.42111}{t}{Det er avgrensa av naboskapet til dei alt busette regionane og av barrierane mot dei forbodne\footnote{Migrering til ny funksjonalitet er asymmetrisk avgrensa av eldre, hardkoda avhengigheiter.}.}
\prop{1.42112}{o}{Naboskapet til dei busette trekkjer til seg; dei forbodne støyter ifrå\footnote{Brukarvennlegheit virkar som ein attraktor, kognitiv overbelasting fungerer som systemets frastøytande horisont.}.}
\prop{1.421121}{o}{Det realiserte rommet er ein funksjon av båe to\footnote{Likevekta mellom eit systems latens og informasjonsgjennomstrøyming stabiliserer arkitekturen.}.}
\prop{1.43}{d}{Dei forbodne regionane er dei som vert utelukka av materielle, strukturelle eller energetiske restriksjonar under gjeldande vilkår\footnote{Konkurransen om den endelege tilstanden i eit transaksjonsminne definerer dei fysisk uframkallelege asymmetriane.}.}
\prop{1.431}{t}{Det som er forbode under eitt sett vilkår, treng ikkje vere forbode under eit anna\footnote{Introduksjon av asynkrone trådar tilgjengeliggjer augneblikkeleg den synkront uoppnåelege arkitekturen.}.}
\prop{1.4311}{o}{Grensa mellom det opne og det forbodne er dynamisk\footnote{Ei oppdatering av den underliggande kompilatoren forskyv heile randsona for handlingsrommet.}.}
\prop{1.43111}{o}{Eit materiale som gjer ein region forboden, kan skiftast ut og soleis opne opp regionen\footnote{Eit paradigmeskifte i lagringsteknologi inverterer kva operasjonar som er dyre og kva som kjem gratis.}.}
\prop{1.43112}{o}{Grensa er ikkje ein eigenskap ved forma; ho er ein eigenskap ved vilkåra\footnote{Nettverkstopologien sin restriksjon ligg i protokollane og maskinvaren, ikkje i datasettet sjølv.}.}